\documentclass[conference]{IEEEtran}

% ── Encoding ─────────────────────────────────────────────────────────────────
\usepackage[utf8]{inputenc}
\usepackage[T1]{fontenc}

% ── Math ─────────────────────────────────────────────────────────────────────
\usepackage{amsmath,amssymb,bm}

% ── Tables ───────────────────────────────────────────────────────────────────
\usepackage{booktabs}
\usepackage{multirow}
\usepackage{array}

% ── Figures ──────────────────────────────────────────────────────────────────
\usepackage{graphicx}
\usepackage{float}

% ── Colors & links ───────────────────────────────────────────────────────────
\usepackage[dvipsnames]{xcolor}
\usepackage{hyperref}
\hypersetup{colorlinks=true,linkcolor=black,citecolor=black,urlcolor=black}

% ── Misc ─────────────────────────────────────────────────────────────────────
\usepackage{microtype}
\usepackage{cite}
\usepackage{algorithm}
\usepackage{algpseudocode}
\usepackage{enumitem}

% ── Draft helpers (remove before submission) ─────────────────────────────────
\newcommand{\TBD}{\textcolor{red}{\textbf{[TBD]}}}

% ── Shorthands ───────────────────────────────────────────────────────────────
\newcommand{\ie}{\textit{i.e.}}
\newcommand{\eg}{\textit{e.g.}}
\newcommand{\etal}{\textit{et~al.}}
\newcommand{\Lin}{L_{\mathrm{in}}}
\newcommand{\Lout}{L_{\mathrm{out}}}

% ─────────────────────────────────────────────────────────────────────────────
\begin{document}

\title{Adaptive Ensemble Learning for Multi-Step Heart Rate Forecasting:\\
       Deep Learning, Foundation Models, and Artificial Immune Systems}

\author{%
  \IEEEauthorblockN{[Author Names]}\\
  \IEEEauthorblockA{[Institution / Department]\\
  [City, Country]\\
  adbyb.es@gmail.com}
}

\maketitle

% ─────────────────────────────────────────────────────────────────────────────
\begin{abstract}
Accurate multi-step heart rate (HR) forecasting during physical activity is
essential for real-time sports performance monitoring and health risk
assessment. We propose a comprehensive ensemble framework that combines six
locally trained deep learning architectures---Temporal Convolutional Network
(TCN), N-BEATS, LSTM, TiDE, Encoder-Decoder LSTM, and iTransformer---with two
fine-tuned time series foundation models (MOMENT and Moirai). Four ensemble
strategies are compared under a strict anti-leakage protocol: simple averaging,
optimisation-based weighted combination, Ridge regression stacking, and a novel
Artificial Immune System adaptive ensemble (AIS-Ensemble). Experiments are
conducted on a real-world multi-subject, multi-session exercise dataset
comprising 39 subjects across 10 structured training protocols. Results show
that ensemble methods consistently outperform individual models, with Ridge
stacking achieving up to an 11.9\% relative MAPE reduction over the
simple-average baseline. An exhaustive search over all 31 non-empty subsets of
five models confirms that stacking dominates optimisation-based weighting across
all combination sizes. The AIS-Ensemble introduces per-subject adaptation and
out-of-distribution detection; while it does not surpass stacking in aggregate
MAPE, it provides a biologically inspired interpretable mechanism for
context-sensitive combination and reveals subject-level predictability patterns.
\end{abstract}

\begin{IEEEkeywords}
heart rate forecasting, time series, deep learning, foundation models, ensemble
learning, artificial immune system, sports analytics
\end{IEEEkeywords}

% ─────────────────────────────────────────────────────────────────────────────
\section{Introduction}
\label{sec:intro}
% ─────────────────────────────────────────────────────────────────────────────

Heart rate (HR) is one of the most informative physiological signals for
exercise intensity quantification, fatigue monitoring, and cardiovascular health
assessment \cite{achten2003}. The ability to forecast future HR trajectories in
real time --- given a short history of recent measurements --- would enable
proactive training-load control, early detection of overexertion, and
personalised pacing strategies.

Classical approaches based on ARIMA or linear regression with physiological
covariates \cite{box2015} struggle with the nonlinear, non-stationary dynamics
of exercise-induced HR responses. Deep learning architectures --- convolutional,
recurrent, and attention-based --- have substantially improved accuracy on
physiological time series \cite{nakamura2020}. More recently, large-scale time
series foundation models (MOMENT \cite{goswami2024moment}, Moirai
\cite{woo2024moirai}) raise the question of whether domain-agnostic pretraining
can match or surpass specialised local models in the physiological domain.

Independently, ensemble methods have proven consistently superior to individual
models in forecasting competitions \cite{makridakis2018m4}. Yet naive averaging
can be harmful when model quality is highly non-uniform, and the optimal
combination strategy may vary across subjects given the substantial inter-individual
variability in HR dynamics.

\textbf{Contributions.} This paper: (1) provides a systematic comparison of
eight models (six local + two foundation) for multi-step HR forecasting at two
horizon lengths; (2) evaluates four ensemble strategies under a rigorous
anti-leakage protocol that includes an exhaustive search over all 31 subsets of
five base models; (3) introduces an AIS-based adaptive ensemble that adapts
combination weights per subject using immunological memory and detects
out-of-distribution windows via negative selection; and (4) provides empirical
evidence that model heterogeneity and ensemble sophistication both matter, with
stacking emerging as the dominant strategy.

% ─────────────────────────────────────────────────────────────────────────────
\section{Related Work}
\label{sec:related}
% ─────────────────────────────────────────────────────────────────────────────

\subsection{Heart Rate Prediction with Deep Learning}

Early physiological signal forecasting relied on ARIMA and linear regression
models \cite{box2015}. Deep learning approaches --- particularly LSTM
\cite{nakamura2020} and convolutional networks --- substantially improved
accuracy on exercise HR by capturing nonlinear temporal patterns. Attention
mechanisms further boosted performance on longer sequences \cite{liu2024itransformer}.
Multi-step (sequence-to-sequence) forecasting remains more challenging than
single-step prediction, as errors accumulate along the horizon.

\subsection{Time Series Foundation Models}

MOMENT \cite{goswami2024moment} and Moirai \cite{woo2024moirai} are
transformer-based models pretrained on diverse large-scale corpora, offering
zero-shot or fine-tuned forecasting. Their advantage in specialised domains
such as physiological signals is unclear: pretraining may not cover
exercise-induced HR dynamics, but the learned representations can still
provide a useful inductive bias, especially when per-subject data is scarce.

\subsection{Ensemble Methods in Forecasting}

Ensemble diversity is the primary driver of improvement over individual models
\cite{makridakis2018m4}. Stacking \cite{wolpert1992} trains a meta-learner on
held-out validation predictions, enabling horizon-dependent combination.
Mixture-of-experts (MoE) \cite{jacobs1991} and algorithm-selection
\cite{rice1976} approaches select or weight models based on input features,
motivating our AIS-Ensemble.

\subsection{Artificial Immune Systems}

AIS are biologically inspired computational frameworks \cite{dasgupta1999}.
The aiNet algorithm \cite{decastro2000} constructs a compact network of memory
cells (antibodies) via clonal selection and hypermutation, learning a compressed
representation of the input distribution. Applications span clustering, anomaly
detection, and classification, but their use as an adaptive ensemble mechanism
for physiological time series forecasting has not been previously reported.

% ─────────────────────────────────────────────────────────────────────────────
\section{Dataset and Preprocessing}
\label{sec:data}
% ─────────────────────────────────────────────────────────────────────────────

\subsection{Data Collection}

The dataset comprises HR time series from \textbf{39 subjects} across
\textbf{10 structured exercise sessions} (Sessions 01--10), including FTP ramp
tests, base conditioning (BC) exercises, and HIIT protocols. Each series
contains 1,841 timesteps sampled at a uniform frequency (HR in bpm), yielding
132 total series prior to splitting.

\subsection{Preprocessing}

Each series is independently standardised (zero mean, unit variance) using
parameters estimated exclusively from the training partition. Stored
standardisation parameters are used to restore predictions to bpm before
computing metrics.

\subsection{Data Partitioning}

The dataset is split chronologically into a \textbf{70\% training} partition
(91 series) used exclusively to train base models, and a \textbf{30\%
evaluation} partition (39 series). The evaluation partition is further divided
into two equal halves:
\begin{itemize}[noitemsep, topsep=2pt]
  \item \textit{ens\_fit} (${\approx}$24,301 windows): used to fit ensemble
        weights, the stacking meta-learner, and AIS memory matrices.
  \item \textit{held\_out} (${\approx}$24,302 windows): used exclusively for
        final metric reporting.
\end{itemize}

Sliding-window instances are generated with equal input and output lengths
$\Lin = \Lout \in \{45, 200\}$ and stride 1. All individual model predictions
are cached before any ensemble fitting, ensuring no ensemble method accesses
held-out targets during training.

% ─────────────────────────────────────────────────────────────────────────────
\section{Models}
\label{sec:models}
% ─────────────────────────────────────────────────────────────────────────────

\subsection{Individual Forecasting Models}

All six local models follow a sequence-to-sequence paradigm: given
$\Lin$ standardised HR values, predict the next $\Lout$ values. Training uses
only the 70\% partition with inner 70/10/20 train/val/test splits.

\textbf{TCN} \cite{bai2018} uses dilated causal convolutions with dilation
factors $\{1,2,4,8\}$, kernel size 12, 32 filters per layer, and residual
connections. Its fully parallelisable training and stable gradients make it
well-suited for HR sequences.

\textbf{N-BEATS} \cite{oreshkin2020} decomposes time series into trend and
seasonality components via stacked fully connected blocks with learned basis
expansions.

\textbf{LSTM} uses two stacked layers with gating mechanisms that capture
long-range dependencies characteristic of HR recovery curves.

\textbf{TiDE} \cite{das2024tide} encodes input history into a compact dense
vector and decodes it into the forecast horizon using MLP layers.

\textbf{EncDec} encodes the input context with one LSTM into a fixed-length
state, then decodes it autoregressively with a second LSTM.

\textbf{iTransformer} \cite{liu2024itransformer} applies self-attention over
the variate dimension rather than the time dimension. In the univariate setting,
this reduces to standard attention over the feature representation.

\subsection{Foundation Models}

\textbf{MOMENT} (\texttt{AutonLab/MOMENT-1-large}, ${\approx}$300M parameters)
\cite{goswami2024moment} is fine-tuned on the 70\% training partition
(context length 512, 5 epochs, $\mathrm{lr}=5{\times}10^{-4}$, AdamW,
forecasting head only).

\textbf{Moirai} (\texttt{moirai-1.1-R-small}) \cite{woo2024moirai} is
fine-tuned with patch size 32 and frequency token \texttt{`h'}; the median of
20 Monte Carlo samples is used as the point forecast.

% ─────────────────────────────────────────────────────────────────────────────
\section{Ensemble Methods}
\label{sec:ensemble}
% ─────────────────────────────────────────────────────────────────────────────

\subsection{Simple Average}

\begin{equation}
  \hat{y}_{\mathrm{avg}}(t) = \frac{1}{N}\sum_{i=1}^{N}\hat{y}_i(t).
\end{equation}

Equal weights serve as a variance-reduction baseline requiring no fitting.

\subsection{Optimisation-Based Weighted Average}

\begin{equation}
  \hat{y}_{\mathrm{opt}}(t) = \sum_{i=1}^{N}w_i\hat{y}_i(t),\quad
  w_i\ge0,\;\sum_i w_i=1,
\end{equation}

where $\mathbf{w}$ minimises MAPE on \textit{ens\_fit} via Nelder-Mead
\cite{nelder1965}. Weights are projected onto the probability simplex after
each update. The optimiser is free to zero out underperforming models.

\subsection{Ridge Regression Stacking}

A Ridge meta-learner ($\alpha=1.0$) is trained on the
$N{\times}\Lout$-dimensional concatenation of base model predictions over
\textit{ens\_fit}:
\begin{equation}
  \hat{y}_{\mathrm{stack}}(t) =
    \mathbf{W}^{\!\top}[\hat{y}_1;\ldots;\hat{y}_N]+\mathbf{b}.
\end{equation}

Unlike static weights, stacking learns horizon-dependent combination, \ie\
it can favour model $i$ at early forecast steps and model $j$ at later steps.

\subsection{AIS-Adaptive Ensemble}
\label{ssec:ais}

The AIS-Ensemble adapts combination weights per context window and per subject.
It has five components:

\textbf{Feature extraction.} A 20-dimensional vector $\boldsymbol{\phi}(x)$ is
extracted from each input window, covering basic statistics (mean, std, min,
max, range), shape (skewness, kurtosis, coefficient of variation), trend (slope,
curvature), autocorrelation at lags $\{1,5,10,20\}$, spectral features
(dominant frequency, spectral entropy), activity (zero-crossing rate, peak
density), and percentiles (Q25, Q75).

\textbf{aiNet clustering.} The aiNet algorithm \cite{decastro2000} learns
$K=30$ memory antibodies $\{c_k\}$ in feature space from \textit{ens\_fit}
windows via clonal selection (5 clones per antibody), somatic hypermutation,
and suppression (threshold $\sigma_s=0.3$) over 50 iterations.

\textbf{Affinity computation.} For a test window with features
$\boldsymbol{\phi}(x)$:
\begin{equation}
  a_k = \exp\!\left(-\frac{\|\boldsymbol{\phi}(x)-c_k\|^2}{2\sigma^2}\right),
  \quad \sigma=1.0.
\end{equation}

\textbf{Per-subject memory calibration.} For each subject $i$, a memory
matrix $M_i\in\mathbb{R}^{K\times N}$ is built on \textit{ens\_fit} by
recording which base model predicted best for each antibody cluster.
$M_i[k,m]$ is normalised to a probability distribution over models.

\textbf{Adaptive inference.} For a window from subject $i$:
\begin{equation}
  \mathbf{w}^*(i,x)=
    \operatorname{softmax}\!\Bigl(\sum_{k=1}^{K}a_k\cdot M_i[k,:]\Bigr),
  \quad
  \hat{y}_{\mathrm{AIS}}=\sum_m w^*_m\hat{y}_m.
\end{equation}

\textbf{Negative selection.} A KD-tree on \textit{ens\_fit} feature vectors
flags test windows beyond the 99th percentile training distance as
out-of-distribution (nonself). These windows fall back to TCN (best individual
model) rather than the adaptive combination.

% ─────────────────────────────────────────────────────────────────────────────
\section{Experimental Setup}
\label{sec:setup}
% ─────────────────────────────────────────────────────────────────────────────

\subsection{Evaluation Metrics}

All metrics are computed on \textit{held\_out} after restoring predictions to
bpm scale:
\begin{itemize}[noitemsep, topsep=2pt]
  \item \textbf{MAPE}: mean absolute percentage error (\%), mean across windows.
  \item \textbf{MAPE\textsubscript{med}}: median MAPE, robust to outlier windows.
  \item \textbf{DTW}: Dynamic Time Warping distance (Sakoe-Chiba $w=20$), mean
        across windows.
  \item \textbf{Pearson}: mean Pearson correlation coefficient across windows.
  \item \textbf{Pearson\textsubscript{med}}: median Pearson correlation.
\end{itemize}

\subsection{Configurations}

Two window configurations are evaluated (Table~\ref{tab:configs}).

\begin{table}[H]
  \caption{Experimental window configurations.}
  \label{tab:configs}
  \centering
  \begin{tabular}{ccc}
    \toprule
    $\Lin=\Lout$ & Ensemble methods & Approx.\ held-out windows \\
    \midrule
    45  & Avg, Opt, Stack        & 24,302 \\
    200 & Avg, Opt, Stack, AIS   & 24,302 \\
    \bottomrule
  \end{tabular}
\end{table}

Additionally, a \textbf{permutation search} with $L=200$ evaluates all 31
non-empty subsets of five models (TCN, N-BEATS, LSTM, TiDE, EncDec;
iTransformer excluded due to poor individual performance) under both
optimisation and stacking.

% ─────────────────────────────────────────────────────────────────────────────
\section{Results}
\label{sec:results}
% ─────────────────────────────────────────────────────────────────────────────

\subsection{Individual Models ($L=45$)}

Table~\ref{tab:individual_45} reports held-out results for individual models
with $\Lin=\Lout=45$.

\begin{table}[ht]
  \caption{Individual model performance, $L=45$, held-out split.}
  \label{tab:individual_45}
  \centering
  \resizebox{\columnwidth}{!}{%
  \begin{tabular}{lccccc}
    \toprule
    Model & MAPE$\downarrow$ & MAPE\textsubscript{med}$\downarrow$ &
            DTW$\downarrow$ & Pearson$\uparrow$ &
            P\textsubscript{med}$\uparrow$ \\
    \midrule
    \textbf{TCN}  & \textbf{3.846} & \textbf{2.881} & \textbf{33.54}
                  & 0.498 & \textbf{0.854} \\
    EncDec        & 3.893 & 2.920 & 34.29 & 0.486 & 0.839 \\
    LSTM          & 3.953 & 2.943 & 34.09 & 0.489 & 0.833 \\
    N-BEATS       & 4.920 & 3.931 & 43.14 & 0.432 & 0.792 \\
    TiDE          & 5.750 & 3.881 & 48.37 & 0.469 & 0.786 \\
    iTransformer  & 8.528 & 6.958 & 77.40 & 0.258 & 0.692 \\
    \midrule
    MOMENT        & \TBD  & \TBD  & \TBD  & \TBD  & \TBD  \\
    Moirai        & \TBD  & \TBD  & \TBD  & \TBD  & \TBD  \\
    \bottomrule
  \end{tabular}}
\end{table}

TCN is the best individual model across all metrics except mean Pearson, where
EncDec, LSTM, and the three ensemble strategies surpass it
(Table~\ref{tab:ensemble_45}). The iTransformer's MAPE (8.53\%) is $2.2\times$
that of TCN (3.85\%), confirming that its variate-dimension attention provides
no benefit in the univariate setting.

\subsection{Ensemble Performance ($L=45$)}

\begin{table}[ht]
  \caption{Ensemble performance, $L=45$, held-out split.}
  \label{tab:ensemble_45}
  \centering
  \resizebox{\columnwidth}{!}{%
  \begin{tabular}{lccccc}
    \toprule
    Method & MAPE$\downarrow$ & MAPE\textsubscript{med}$\downarrow$ &
             DTW$\downarrow$ & Pearson$\uparrow$ &
             P\textsubscript{med}$\uparrow$ \\
    \midrule
    Simple Average         & 4.352 & 3.421 & 38.00 & 0.502 & 0.857 \\
    Weighted (Opt.)        & 3.801 & 2.851 & 33.37 & \textbf{0.503} & 0.856 \\
    \textbf{Stacking (Ridge)} & \textbf{3.754} & 2.906 & \textbf{33.27}
                              & 0.496 & 0.833 \\
    \bottomrule
  \end{tabular}}
\end{table}

Simple averaging degrades performance relative to the best individual model
(MAPE 4.35\% vs.\ TCN 3.85\%), because uniformly averaging over low-quality
models such as iTransformer dilutes the signal. The optimised weights recover
this by assigning zero weight to TiDE and iTransformer (learned weights:
TCN=0.476, EncDec=0.323, LSTM=0.151, N-BEATS=0.050, TiDE=0, iTransformer=0).
Ridge stacking achieves the lowest MAPE (3.754\%), an \textbf{11.9\% relative
improvement} over simple averaging, by learning horizon-dependent combination
weights.

A slight trade-off is observed: stacking reduces amplitude errors (MAPE, DTW)
at a marginal cost to Pearson\textsubscript{med} (0.833 vs.\ 0.857 for
averaging), suggesting that the meta-learner slightly distorts trajectory shape
in some windows.

\subsection{Individual Models and Ensembles ($L=200$)}

\begin{table}[ht]
  \caption{Individual model performance, $L=200$, held-out split.}
  \label{tab:individual_200}
  \centering
  \resizebox{\columnwidth}{!}{%
  \begin{tabular}{lccccc}
    \toprule
    Model & MAPE$\downarrow$ & MAPE\textsubscript{med}$\downarrow$ &
            DTW$\downarrow$ & Pearson$\uparrow$ &
            P\textsubscript{med}$\uparrow$ \\
    \midrule
    TCN          & \TBD & \TBD & \TBD & \TBD & \TBD \\
    EncDec       & \TBD & \TBD & \TBD & \TBD & \TBD \\
    LSTM         & \TBD & \TBD & \TBD & \TBD & \TBD \\
    N-BEATS      & \TBD & \TBD & \TBD & \TBD & \TBD \\
    TiDE         & \TBD & \TBD & \TBD & \TBD & \TBD \\
    iTransformer & \TBD & \TBD & \TBD & \TBD & \TBD \\
    MOMENT       & \TBD & \TBD & \TBD & \TBD & \TBD \\
    Moirai       & \TBD & \TBD & \TBD & \TBD & \TBD \\
    \bottomrule
  \end{tabular}}
\end{table}

\begin{table}[ht]
  \caption{Ensemble performance, $L=200$, held-out split.}
  \label{tab:ensemble_200}
  \centering
  \resizebox{\columnwidth}{!}{%
  \begin{tabular}{lccccc}
    \toprule
    Method & MAPE$\downarrow$ & MAPE\textsubscript{med}$\downarrow$ &
             DTW$\downarrow$ & Pearson$\uparrow$ &
             P\textsubscript{med}$\uparrow$ \\
    \midrule
    Simple Average       & 5.383 & \TBD & \TBD & \TBD & \TBD \\
    Weighted (Opt.)      & 4.977 & \TBD & \TBD & \TBD & \TBD \\
    Stacking (Ridge)     & 4.221 & \TBD & \TBD & \TBD & \TBD \\
    \textbf{AIS-Ensemble}& \TBD  & \TBD & \TBD & \TBD & \TBD \\
    \bottomrule
  \end{tabular}}
\end{table}

At $L=200$, all methods show higher MAPE than at $L=45$, consistent with
increasing uncertainty at longer horizons. Stacking (MAPE 4.221) maintains its
advantage over optimisation-based weighting (4.977) and averaging (5.383). The
absolute gap between stacking and simple averaging is larger at $L=200$
(1.16 pp) than at $L=45$ (0.60 pp), suggesting that horizon-dependent
combination is more valuable for longer forecasts.

\subsection{Permutation Search ($L=200$)}

Table~\ref{tab:permutations} summarises the top-5 combinations for each
ensemble method from the exhaustive search over all 31 non-empty subsets of
five models.

\begin{table}[ht]
  \caption{Top-5 combinations from exhaustive search, $L=200$. Best MAPE in
           bold. iTransformer excluded from search.}
  \label{tab:permutations}
  \centering
  \resizebox{\columnwidth}{!}{%
  \begin{tabular}{lcc}
    \toprule
    Combination & MAPE\textsubscript{Opt}$\downarrow$ &
                  MAPE\textsubscript{Stack}$\downarrow$ \\
    \midrule
    N-BEATS + TiDE                         & 5.006 & \textbf{4.627} \\
    TCN + TiDE                             & 4.985 & 4.759 \\
    LSTM + TiDE                            & 5.022 & 4.758 \\
    N-BEATS + LSTM                         & 5.014 & 4.877 \\
    TCN + N-BEATS                          & 5.022 & 4.911 \\
    \midrule
    TCN+N-BEATS+LSTM+TiDE+EncDec (all 5)  & 4.977 & 4.221 \\
    \bottomrule
  \end{tabular}}
\end{table}

Several findings emerge. First, stacking consistently dominates
optimisation-based weighting across all 31 combinations --- a systematic finding
that validates the advantage of horizon-aware meta-learning. Second, the
best 2-model stacking combination (N-BEATS + TiDE, MAPE 4.627) outperforms the
best 5-model optimisation combination (MAPE 4.977), illustrating that ensemble
\emph{quality} matters more than ensemble \emph{size} under optimisation.
Third, TiDE --- which performed poorly as an individual model (MAPE 5.75 at
$L=45$) --- appears frequently in top stacking pairs, suggesting that its
diverse error pattern is exploited by the meta-learner even when its standalone
accuracy is low.

\subsection{AIS-Ensemble Analysis}

The AIS-Ensemble operates at the intersection of ensemble learning and
immunological computation. \TBD (full quantitative results pending
re-execution of \texttt{ensemble\_ais.ipynb}.) Preliminary diagnostics reveal:

\textbf{Out-of-distribution detection.} Approximately 10--20\% of test windows
are classified as nonself (out-of-distribution) per subject, indicating that
the exercise patterns in \textit{held\_out} are not always well-represented by
the \textit{ens\_fit} calibration distribution. These windows fall back to TCN.

\textbf{Per-subject weight variation.} The subject memory matrices $M_i$ show
diverse distributions across subjects: some subjects have memory matrices
concentrated on TCN (uniform, steady-state HR profiles), while others distribute
weight across TCN and EncDec (complex recovery patterns).

\textbf{Aggregate performance.} Although the AIS-Ensemble does not improve
upon stacking in aggregate MAPE --- consistent with the finding that
per-window adaptation is not the primary bottleneck in this dataset --- it
produces subject-level insights not available from static ensemble methods. In
particular, the nonself ratio per subject correlates with inter-subject
variability in individual model accuracy, suggesting that OOD detection could
serve as a proxy for subject-specific model reliability.

\subsection{Horizon Comparison}

Comparing $L=45$ and $L=200$ results: (i) MAPE increases with horizon for all
methods; (ii) the relative advantage of stacking over simple averaging grows
with horizon (11.9\% at $L=45$ vs.\ \TBD at $L=200$); (iii) the relative
ranking of individual models is consistent across both horizons, with TCN,
EncDec, and LSTM forming the top tier.

\subsection{Foundation Models}

\TBD\ (pending fine-tuning notebook re-execution). We hypothesise that:
(i) zero-shot foundation model performance will trail well-trained local models
due to domain mismatch; (ii) fine-tuning will substantially close the gap;
(iii) foundation models may provide the greatest benefit for subjects with the
fewest training sessions, where the pretrained representation compensates for
data scarcity.

% ─────────────────────────────────────────────────────────────────────────────
\section{Discussion}
\label{sec:discussion}
% ─────────────────────────────────────────────────────────────────────────────

\textbf{Stacking vs.\ optimisation.} The systematic superiority of Ridge
stacking across all model subsets and both horizon lengths is a central finding.
The key advantage is horizon-awareness: the Ridge meta-learner can assign
different effective weights to each model at each forecast step. For HR
forecasting, this is meaningful because different architectures capture
different temporal scales --- TCN excels at short-range dynamics via dilated
convolutions, while EncDec captures longer response curves via its latent-space
encoding.

\textbf{Danger of simple averaging with heterogeneous quality.} The degradation
of simple averaging relative to the best individual model (at $L=45$) is a
practical warning: when ensemble members span a wide quality range (TCN MAPE
3.85\% vs.\ iTransformer 8.53\%), equal-weight combination shifts the
prediction towards the low-quality tail. Practitioners should use at least
optimised weights or stacking when combining heterogeneous forecasters.

\textbf{AIS-Ensemble as a diagnostic tool.} While the AIS-Ensemble does not
improve aggregate MAPE beyond stacking, it offers a unique capability:
subject-level adaptation and uncertainty quantification via out-of-distribution
detection. In a real deployment scenario, the nonself ratio per session could
trigger an alert (e.g., flagging unusual HR responses), complementing the
quantitative forecast.

\textbf{Short vs.\ long horizon.} The consistent superiority of $L=45$ over
$L=200$ in absolute MAPE (3.75\% vs.\ 4.22\% for stacking) reflects the
fundamental difficulty of long-horizon physiological forecasting. Shorter
windows may also benefit from more training instances per series (higher window
density), which can improve all downstream models and ensemble weights.

% ─────────────────────────────────────────────────────────────────────────────
\section{Conclusion}
\label{sec:conclusion}
% ─────────────────────────────────────────────────────────────────────────────

We presented a systematic study of individual deep learning and foundation
models with four ensemble strategies for multi-step HR forecasting. Key
conclusions are:

\begin{enumerate}[noitemsep, topsep=2pt]
  \item TCN is the strongest individual model; iTransformer and TiDE are
        significantly weaker in the univariate setting.
  \item Ridge stacking consistently outperforms optimisation-based weighting
        and simple averaging, achieving an 11.9\% relative MAPE reduction over
        the simple-average baseline at $L=45$.
  \item Simple averaging degrades performance when model quality is
        non-uniform --- practitioners must use weighted or stacking methods.
  \item The permutation search confirms stacking dominance across all 31 model
        subsets; two-model stacking can outperform five-model optimisation.
  \item The AIS-Ensemble offers subject-level interpretability and
        out-of-distribution detection, complementing aggregate accuracy with
        patient-specific diagnostics.
\end{enumerate}

Future work will incorporate physiological covariates (power, cadence, RPE) as
ensemble conditioning signals, explore online AIS memory updates as sessions
accumulate, and extend the framework to probabilistic HR forecasting.

% ─────────────────────────────────────────────────────────────────────────────
\appendix
% ─────────────────────────────────────────────────────────────────────────────

\section{Model Hyperparameters}

\begin{table}[H]
  \caption{Local model hyperparameters.}
  \label{tab:hparams}
  \centering
  \resizebox{\columnwidth}{!}{%
  \begin{tabular}{ll}
    \toprule
    Model & Key hyperparameters \\
    \midrule
    TCN          & filters=32, kernel\_size=12, dilations=\{1,2,4,8\} \\
    N-BEATS      & stacks=2, blocks=3, layers=4, width=256 \\
    LSTM         & units=[128,64], dropout=0.2 \\
    TiDE         & enc\_dim=256, dec\_dim=256, temporal\_width=4 \\
    EncDec       & enc\_units=128, dec\_units=128, latent=64 \\
    iTransformer & d\_model=128, heads=8, d\_ff=256, layers=2 \\
    MOMENT       & context=512, epochs=5, lr=$5{\times}10^{-4}$ \\
    Moirai       & patch\_size=32, samples=20 \\
    \bottomrule
  \end{tabular}}
\end{table}

\section{AIS Hyperparameters}

\begin{table}[H]
  \caption{AIS-Ensemble hyperparameters.}
  \label{tab:ais_hparams}
  \centering
  \begin{tabular}{ll}
    \toprule
    Parameter & Value \\
    \midrule
    Antibodies ($K$)         & 30   \\
    Feature dimension        & 20   \\
    Affinity $\sigma$        & 1.0  \\
    Suppression $\sigma_s$   & 0.3  \\
    Clones / antibody        & 5    \\
    Max iterations           & 50   \\
    OOD threshold            & 99th pct. \\
    Fallback model           & TCN  \\
    \bottomrule
  \end{tabular}
\end{table}

% ─────────────────────────────────────────────────────────────────────────────
\bibliographystyle{IEEEtran}
\bibliography{references}
% ─────────────────────────────────────────────────────────────────────────────

\end{document}
